\section{Related Work}
\label{sec:related}

Section~\ref{sec:basepaper} covered the works classified in the base survey. This section looks at work beyond the survey, grouped by the problem each line of work solves, and positions TrustTrail against it. Table~\ref{tab:related} summarises the comparison.

\subsection{Privacy-preserving single sign-on}
A standard OIDC provider learns every service that a user logs in to, which lets the IdP build a profile of the user's activity. Several designs remove this. SPRESSO by Fett, Küsters and Schmitz is a web SSO system in which the IdP does not learn which relying party the user logs in to~\cite{fett2015spresso}. Hammann, Sasse and Basin propose a privacy-preserving extension of OIDC with the same aim~\cite{hammann2020privacy}, and Kroschewski and Lehmann address how a relying party can be authenticated in OIDC without revealing it to the IdP~\cite{kroschewski2023implicit}. UPPRESSO, so far available as a preprint, aims at SSO that is both untraceable by the IdP and unlinkable across relying parties~\cite{guo2021uppresso}. On the security side, Fett et al.\ gave a formal analysis of OAuth~2.0~\cite{fett2016oauth}, Mainka et al.\ systematised attacks on OIDC~\cite{mainka2017sok}, and Wang, Chen and Wang showed through traffic analysis that deployed SSO systems had serious logic flaws~\cite{wang2012signing}.

These works protect the user from the IdP, or protect the protocol from attackers. TrustTrail addresses a different question: what the relying parties receive, under which promise, and whether that can be checked later. The two concerns are complementary. In a campus federation the IdP is the institute itself and already holds the student's records, so IdP tracking is less of a concern than over-release to the many services around it. We list IdP tracking as a non-goal in Section~\ref{sec:threat}; a deployment that needs both could combine TrustTrail with one of these protocols.

\subsection{Anonymous credentials and self-sovereign identity}
Chaum proposed in 1985 that individuals use a different pseudonym with each organisation and prove facts about themselves with credentials that cannot be linked across organisations~\cite{chaum1985security}. Camenisch and Lysyanskaya gave an efficient anonymous credential system~\cite{camenisch2001efficient}, which IBM implemented as Idemix~\cite{camenisch2002idemix}; Microsoft's U-Prove offers similar selective disclosure~\cite{paquin2013uprove}. The PRIME, PrimeLife and ABC4Trust projects in the base survey built on these systems. More recently, self-sovereign identity places credentials in a wallet controlled by the user~\cite{allen2016path,muhle2018survey}, and the W3C Verifiable Credentials data model standardises their format~\cite{w3cvc2}.

Credential systems give strong unlinkability and fine-grained disclosure, often stronger than pairwise identifiers, because even the issuer cannot link presentations. They also move work to the user: the student needs a wallet, and every service needs a verifier. More importantly for our goals, a credential presentation leaves no record that the user can later show to a third party, and nothing ties the verifier's use of the data to a declared purpose. TrustTrail takes the opposite trade-off. It keeps standard OIDC on the service side and gives up some unlinkability (the institute's IdP and TrustTrail know the mapping between pseudonyms), in exchange for a central point where purpose, retention and the release record can be enforced and kept.

\subsection{Policy-based control of disclosure}
A large part of the base survey concerns policies that state how released data may be used. Sticky policies attach machine-readable usage conditions to data so that the conditions travel with it~\cite{pearson2011sticky}. Chadwick and Fatema built a privacy-preserving authorisation system in which users' policies are stuck to their data and enforced by the cloud provider's policy decision points~\cite{chadwick2012privacy}. Orawiwattanakul et al.\ let users filter which attributes a Shibboleth IdP releases~\cite{orawiwattanakul2010user}, and Suriadi, Foo and Jøsang designed a user-centric federated SSO system with privacy requirements on the attribute exchange~\cite{suriadi2009usercentric}. Sánchez et al.\ combined privacy preferences with dynamic federation for consumer cloud services~\cite{sanchez2012enhancing}.

In these systems, the user or the IdP states the policy, and the service is trusted to honour it once it has the data, unless trusted hardware enforces it. TrustTrail inverts the direction of the commitment: the \emph{service} states its purpose and retention, signs the statement with its own key, and TrustTrail checks every release against it. Usage control generalises these ideas. The UCON\textsubscript{ABC} model of Park and Sandhu adds obligations and conditions that are checked before and during use, and can be re-evaluated while access continues~\cite{park2004ucon}. Enforcing such obligations requires a reference monitor wherever the data is used. TrustTrail has a reference monitor only at the release point. Neither approach controls the data once it has been released. TrustTrail adds evidence: the service's signature makes the promise attributable, and the ledger records which data was released under which promise.

\subsection{Consent and receipts}
The Kantara Initiative's Consent Receipt specification proposes that a user receive a record, in a standard format, every time they give consent, listing the controller, the purposes and the data categories~\cite{kantarareceipt}. TrustTrail's release receipts follow the same idea but are produced per release, signed by the release service, and linked to a ledger entry and, through anchoring, to a public chain. Research on consent interfaces explains why the record matters. Utz et al.\ showed in a field study that the design of GDPR consent notices strongly changes what users accept~\cite{utz2019uninformed}, and Felt et al.\ found that few smartphone users pay attention to, or correctly understand, permission prompts~\cite{felt2012attention}. Applications also tend to request more permissions than they use~\cite{felt2011permissions}, and OAuth is often misused in mobile applications~\cite{chen2014oauth}. A consent screen alone is therefore a weak control. TrustTrail's design does not rely on the student reading a prompt: the profile cap and the declaration are checked whether or not the student pays attention, and the dashboard and receipts give a record that can be read later.

\subsection{Tamper-evident logging}
Secure audit logs have a long history. Schneier and Kelsey used hash chains and evolving keys so that log entries written before a compromise cannot be changed undetectably~\cite{schneier1999secure}. Ma and Tsudik proposed forward-secure sequential aggregate signatures for the same purpose, with compact verification~\cite{ma2007forward}. Crosby and Wallach designed history trees, in which the logger commits to the whole log after each append and auditors can check both membership of an entry and consistency between two commitments efficiently~\cite{crosby2009efficient}. Certificate Transparency applies Merkle-tree logs at Internet scale, with monitors and auditors that check that a log only grows~\cite{laurie2014certificate,rfc6962,rfc9162}.

These systems assume an outside party that receives and keeps commitments, such as an auditor, a monitor or a gossip protocol. TrustTrail has two such parties: students, who hold receipts that contain the hash of their own entries, and a public blockchain, which holds the Merkle root of each batch. It uses the RFC~6962 leaf and node hashing and leaves consistency proofs between successive trees for future work, because the contract already refuses any batch that does not continue exactly where the previous one ended.

A log's head can also be fixed without a blockchain. The operator can sign the head hash and publish it, for example on a web page or a mailing list; a trusted timestamping authority can sign it following RFC~3161~\cite{rfc3161}; or independent witnesses can cosign it, as CT monitors do. Each option detects the same rewrites and truncations as anchoring, as long as someone outside the operator's control keeps the published value. They differ in whom the verifier must trust and in cost. Section~\ref{sec:discussion} compares them and explains why we chose a public chain.

\subsection{Blockchains for privacy and accountability}
Zyskind, Nathan and Pentland proposed using a blockchain as an access-control manager for personal data held off chain~\cite{zyskind2015decentralizing}. MedRec uses Ethereum smart contracts to record permissions over medical records stored by providers~\cite{azaria2016medrec}. Neisse, Steri and Nai-Fovino studied blockchain-based data accountability and provenance tracking with the GDPR in mind~\cite{neisse2017blockchain}, and Truong et al.\ proposed a blockchain-based solution for GDPR-compliant personal data management~\cite{truong2020gdpr}. Timestamping services such as OpenTimestamps and Chainpoint anchor Merkle roots of arbitrary documents in public blockchains~\cite{todd2016opentimestamps,chainpoint}.

Many of the blockchain-based designs keep access rights or usage logs themselves on chain. This makes the record verifiable but raises a conflict with data protection law: information written to a public chain cannot be erased, and even pseudonymous log records can be personal data. TrustTrail writes only Merkle roots and entry ranges to the chain. A root reveals nothing about the entries it covers, and the ledger with its pseudonyms stays under the institute's control. This places TrustTrail closer to the anchoring services than to the blockchain access-control designs, with the difference that the anchored data is an accountability ledger with an enforced contiguity rule.

\subsection{Trust and reputation}
Jøsang, Ismail and Boyd survey trust and reputation systems for online services~\cite{josang2007survey}. Most of them aggregate ratings from many users. In the base survey, trust appears as provider reputation shown to the user~\cite{werner2016model} and as static classes of federation membership~\cite{ferdous2015dynamic}. TrustTrail's trust score is simpler and more local: it is computed only from facts that TrustTrail observes directly (registered privacy policy, registration metadata and recorded violations), and its only effect is to lower the most revealing profile the service can obtain. The weights are starting values, which Section~\ref{sec:discussion} discusses.

\begin{table}[t]
\centering
\caption{TrustTrail compared with related approaches (our reading)}
\label{tab:related}
\small
\begin{tabularx}{\linewidth}{@{}>{\raggedright\arraybackslash}p{3.0cm} >{\raggedright\arraybackslash}p{2.6cm} >{\raggedright\arraybackslash}X >{\raggedright\arraybackslash}X@{}}
\toprule
Approach & Examples & What it protects & What it leaves open for our goals \\
\midrule
Privacy profiles as OIDC scopes & \cite{werner2015approach,werner2016model} & Less information released, chosen per interaction & Profile requested, not capped; no linkage protection; no record of releases \\
Privacy-preserving SSO & \cite{fett2015spresso,hammann2020privacy,guo2021uppresso} & User's activity hidden from the IdP & What relying parties receive and do; no purpose binding or release record \\
Anonymous credentials, SSI & \cite{camenisch2002idemix,paquin2013uprove,w3cvc2} & Strong unlinkability, even against the issuer; selective disclosure & Needs wallets and verifiers; nothing records the verifier's stated purpose or what it received for later audit \\
Sticky and user policies & \cite{pearson2011sticky,chadwick2012privacy,orawiwattanakul2010user} & Usage conditions attached to data & Enforcement after release; no evidence the service accepted the policy \\
Consent receipts & \cite{kantarareceipt} & User gets a structured, possibly signed record of each consent & One receipt per consent, not per release; no shared log that makes omissions or edits visible \\
Usage control & \cite{park2004ucon} & Obligations checked before and during use & Needs a reference monitor wherever data is used \\
Tamper-evident logs & \cite{schneier1999secure,crosby2009efficient,laurie2014certificate} & Integrity of a log against later edits & Needs external auditors; does not decide what is released \\
Blockchain access control & \cite{zyskind2015decentralizing,azaria2016medrec,truong2020gdpr} & Verifiable permissions and logs & Records or rights on chain (erasure conflict); new client software \\
\textbf{TrustTrail} & this paper & Profile cap, pairwise IDs, signed purpose, tamper-evident release record, trust feedback & IdP sees all logins; use after release; operator controls the pseudonym map \\
\bottomrule
\end{tabularx}
\end{table}
